% Generated draft: requires mathematical and source review.
\documentclass[11pt]{article}
\usepackage[margin=25mm]{geometry}
\usepackage{fontspec}
\setmainfont{Latin Modern Roman}
\usepackage{amsmath,amssymb,mathtools,unicode-math}
\setmathfont{Latin Modern Math}
\usepackage{enumitem,needspace,xurl,hyperref}
\setlength{\parindent}{0pt}
\setlength{\parskip}{5pt}
\setlength{\emergencystretch}{3em}
\newcommand{\N}{\mathbb{N}}
\newcommand{\Z}{\mathbb{Z}}
\newcommand{\Q}{\mathbb{Q}}
\newcommand{\R}{\mathbb{R}}
\newcommand{\C}{\mathbb{C}}
\newcommand{\F}{\mathbb{F}}
\newcommand{\A}{\mathbb{A}}
\newcommand{\PP}{\mathbb{P}}
\newcommand{\E}{\mathbb{E}}
\newcommand{\eps}{\varepsilon}
\newcommand{\dd}{\,\mathrm{d}}
\begin{document}
\section*{UnsolvedMath ID: 30004565}
\textbf{Source problem code:} catalogue entry
\subsection{Definitions and mathematical statement}
The supplied question is: ``For nonlocal self-avoidance energies on embeddings, do minimizers or stationary points exist within prescribed isotopy classes?'' The record does not specify an energy functional, the source and target of the embeddings, their regularity, the meaning of isotopy, or the admissible variations. Consequently, it determines an existence problem schema rather than a fully quantified conjecture.

To express this schema, let \(M\) and \(N\) denote the intended source and ambient spaces, and let \(\mathcal A\) denote the intended admissible space of embeddings \(f\colon M\hookrightarrow N\), with its regularity and topology still to be specified. An embedding is a map that is a homeomorphism onto its image, with any additional differential regularity required by the chosen setting. Once the intended notion of isotopy is fixed, write \([f]_{\mathrm{iso}}\) for the isotopy class of \(f\). For a prescribed class \(\kappa\), put
\[
\mathcal A_\kappa=\{f\in\mathcal A:[f]_{\mathrm{iso}}=\kappa\}.
\]
Neither ambient isotopy nor isotopy through embeddings is selected here, because the record does not distinguish them. Let \(E\) denote the intended nonlocal self-avoidance energy on \(\mathcal A\); its formula and hypotheses remain unspecified.

The minimizer question asks whether there exists \(f_*\in\mathcal A_\kappa\) such that
\[
E(f_*)\leq E(f)\qquad\text{for every }f\in\mathcal A_\kappa.
\]
If the intended formulation requires finite energy, it must additionally require \(E(f_*)\in\R\). The source does not explicitly state this requirement or specify the range of \(E\).

For a differentiable variational formulation, the stationary-point question asks whether there exists \(f_*\in\mathcal A_\kappa\) such that
\[
\left.\frac{\dd}{\dd t}E(f_t)\right|_{t=0}=0
\]
for every admissible differentiable variation \((f_t)_{|t|<\delta}\) through \(f_0=f_*\), with \(\delta>0\), preserving the intended constraints and isotopy class. This formula applies only after the variation class and the existence of the indicated first derivatives have been specified.

The unresolved target is existence of minimizers, stationary points, or both, for the intended energies and prescribed isotopy classes. The record does not determine whether this is asserted for every such energy and class, only for a particular family, or as a request to identify sufficient hypotheses. No universal existence assertion is made here. The title refers to knots, while the supplied statement concerns embeddings without specifying their dimensions; the knot setting is therefore not imposed.
\subsection{Short English statement}
The problem asks whether a nonlocal self-avoidance energy attains its minimum, or has a stationary embedding, within a prescribed isotopy class. A precise conjecture requires the energy, embedding spaces, constraints, and meaning of stationarity to be specified. The supplied partially solved label is dated curation requiring independent review.
\subsection{Sources}
\begin{itemize}
\item[\textnormal{[1]}] Local UnsolvedMath registry, \texttt{data/unsolvedmath/problems.json}, permanent ID \(30004565\), source problem code \texttt{catalogue entry}, title ``Minimizers of Nonlocal Knot Self-Avoidance Energies.'' The supplied record attributes the question to Oberwolfach Reports workshop discussions; no individual proposer is supplied. Its proposed year is \(2020\).
\item[\textnormal{[2]}] Supplied report citation: \emph{Mini-Workshop: Nonlocal Analysis and the Geometry of Embeddings} (2021), Oberwolfach Reports, pp.\ 1823--1856, DOI \texttt{10.4171/owr/2020/37}. \url{https://doi.org/10.4171/owr/2020/37}. This is also listed among the literature sources as ``OWR 2020/37 report.'' The supplied citation year and the record's proposed year are retained without resolving their relationship. No problem-specific page or passage is supplied.
\item[\textnormal{[3]}] Supplied literature source: \emph{The gradient flow of O'Hara's knot energies}. \url{https://doi.org/10.1007/s00208-017-1540-4}.
\item[\textnormal{[4]}] Supplied literature source: \emph{Palais--Smale condition in knot theory}. \url{https://arxiv.org/abs/2505.02719}.
\item[\textnormal{[5]}] Supplied literature assessment, checked at \texttt{2026-08-22}: status ``partially solved.'' The assessment states that existence of minimizers and stationary knots is known for broad classes of O'Hara-type energies and many isotopy classes, while the source's qualification remains valid for some energies or isotopy constraints. This is supplied dated curation derived from the OWR material, not an independently verified assessment. The record expressly requests independent verification of time-sensitive status and cited claims; it does not identify the precise remaining energies or constraints.
\end{itemize}
\end{document}
